\chapter{Introduction}\label{ch:introduction}
\graphicspath{{./figs/Introduction/}{./figs/Aim1/}{./figs/Aim2/}{./figs/Preliminary/}}

\section{Motivation}

Biomimetic robots provide a platform for investigating how mechanical structure and control interact to generate movement, offering complementary insights to simulation-based studies \citep{andrews_general_1985, millard_flexing_2013, arnold_model_2010}. Wheeled robots remain the platform of choice in structured environments, where their speed, efficiency, and payload capacity are difficult to match; in the unstructured terrain of the natural and built world, legged machines hold the advantage, and biomimetic legs that reproduce human force capabilities have the potential to excel where wheels cannot \citep{frey_legged_2026, ha_learning_2025}. Musculoskeletal humanoids and limb systems driven by tendon-like actuators have been shown to support biologically meaningful experiments on motion generation, internal force transmission, and neural control hypotheses \citep{asano_design_2017, asano_musculoskeletal_2019, hitzmann_anthropomorphic_2018, liu_robotic_2018, chen_realizing_2019}. The dual perspective of using robots to learn about biology, and using biology to design more adaptive robots, has been emphasized in both biomimetic robotics and biomechanics \citep{ijspeert_amphibious_2020, asano_design_2017}. With such robotic systems, experiments can be performed that would not be practical or ethical with human test subjects, such as removing a specific muscle from the system or blocking a sensory feedback pathway. Experiments of this kind allow direct testing of how the components of the neuromechanical system interact, and they offer new opportunities for testing neural control theories that are faster and cheaper than human observation studies \citep{shin_constructive_2018, goldsmith_investigating_2021, schilling_adaptive_2022}.

Artificial muscles are a key enabling technology for this class of robot. Among the available actuation technologies, McKibben-style braided pneumatic actuators (BPAs) combine low mass, high force-to-weight ratio, and intrinsic compliance with a force--length curve that qualitatively resembles that of skeletal muscle \citep{liang_comparative_2020, hunt_modeling_2017}. BPAs have been used in quadrupedal robots \citep{aschenbeck_design_2006, hunt_modeling_2017, scharzenberger_design_2019}, musculoskeletal humanoids \citep{asano_design_2017, hitzmann_anthropomorphic_2018}, and dexterous manipulation systems \citep{chen_realizing_2019}. However, BPAs also differ from biological muscle in important ways: their maximum tensile force occurs at resting length rather than at an optimal fiber length, their maximum contraction ratio is substantially lower than that of human muscle, and their force output is a highly nonlinear function of internal pressure, contraction, and loading state \citep{hunt_modeling_2017, sarosi_comparative_2017, liang_comparative_2020}. The author's experience with \qty{10}{\mm} actuators made this gap concrete: manufacturer specifications overstate the force and contraction available at the short resting lengths a leg design requires, published models fitted at one cut length do not transfer to other resting lengths, and designs carried out against those specifications delivered less range of motion and joint torque than intended, forcing iteration on hardware rather than at the design stage. Designing a biomimetic robot around BPAs therefore requires quantitative, predictive models of the actuator and of the joint torques it can produce.

\section{The Biomimetic Humanoid Robot Project}\label{sec:project}

This dissertation is part of a long-running effort at Portland State University to build a biomimetic humanoid robot from the trunk down, actuated by pneumatic artificial muscles and controlled by a synthetic neural network (SNN). The project's guiding hypothesis is that a robot with human-like biomechanics --- human-like muscle numbers, architecture, and routing, human-like joint ranges of motion, and human-like isometric torque profiles --- controlled by a biologically grounded neural architecture, will be capable of biomimetic movement and will serve as a testbed for neuromechanical hypotheses, including loss-of-function experiments that are impossible in living animals. A by-product of this work is expected to be improved prosthetic and orthotic device design.

The present work builds directly on a series of studies by the author and collaborators. In the author's master's project, the kinematics and dynamics of a biomimetic robot leg were matched against the Gait2392 OpenSim model, artificial muscle attachment and wrapping locations were determined for the robot legs, and MATLAB tooling was developed for evaluating muscle lengths, moment arms, and isometric joint torques over the full range of motion \citep{bolen_bipedal_2019, bolen_determination_2019}. That work showed that simply transplanting human muscle attachment locations onto a BPA-actuated robot does not reproduce human torque profiles. Hand placement of the muscle origin and insertion locations for 27 muscle-actuators per side was required. In parallel, a biomimetic four-bar knee linkage with a migrating instantaneous center of rotation was designed and validated by Steele \citep{steele_development_2017, steele_biomimetic_2018, steele_biomimetic_2018-1}, and a canine-inspired quadruped platform (Muscle Mutt, formerly  called DoggyDeux) with antagonistic BPA pairs and synthetic neural network control was developed by \citet{scharzenberger_design_2019}, extending an earlier demonstration by the Agile and Adaptive Robotics Laboratory (AARL) of central pattern generator (CPG) control of a pneumatically actuated robot \citep{hunt_development_2017}.


\section{Problem Statement}\label{sec:problem}

Within this context, three gaps stand between the existing robot design and a robot that can walk under neural control.

First, the actuator force models available at the start of this work were inadequate for design purposes. The manufacturer's simulation tool over-predicts the force of short BPAs by large margins, and the published empirical models are fitted to a single resting length each, so they do not accurately predict how force changes when a designer chooses a different actuator length \citep{sarosi_comparative_2017, martens_modeling_2017, hunt_modeling_2017}. Because muscle resting length is one of the primary degrees of freedom available to the robot designer, this knowledge-gap reduces confidence that a given design will perform as expected once it is built. 

Second, actuator force alone does not determine joint torque. The robot knee applies BPA force through brackets, artificial tendons, and wrapping paths. It was unknown how much these effects reduce usable torque relative to the rigid-body prediction, or how to best correct for them. Additionally, the biomimetic knee's optimized muscle routing \citep{morrow_optimization_2020} had never been validated experimentally or compared against human isometric torque data.

Third, no simulation pipeline existed for closing the loop between a musculoskeletal robot model of this kind and a biologically realistic neural controller. The Neuromechanical Modeling literature provides the ingredients --- two-level spinal CPG architectures \citep{rybak_modelling_2006, mccrea_organization_2008}, neuromechanical simulation frameworks \citep{markin_afferent_2010}, and synthetic nervous system tooling \citep{nourse_sns_2023} --- but they had not been assembled for a BPA-driven bipedal robot with human-like routing.


\section{Research Objectives}\label{sec:objectives}


This dissertation addresses the first two gaps experimentally and develops the tools and framework needed to investigate the third. The objectives are:

\begin{enumerate}
    \item \textbf{Characterize the isometric force of Festo BPAs} with uninflated diameters of \qtylist{10;20}{\mm} over a wide range of resting lengths, contractions, and pressures; derive improved maximum-force and normalized force--strain--pressure models; and compare them against existing models (Chapters~\ref{ch:methods} and~\ref{ch:results}, Section~\ref{sec:results_force}).
    
    \item \textbf{Validate and improve the prediction of isometric knee torque} produced by BPAs on pinned-joint and biomimetic four-bar knees; identify and model losses caused by constant length offsets, series compliance, and actuator wrapping; and compare the resulting torque envelopes with human capabilities (Methods Sections~\ref{sec:robot_model}, \ref{sec:torque_equipment}, \ref{sec:hybrid_method}, \ref{sec:improved_model}, and~\ref{sec:optimization}; Results Section~\ref{sec:results_torque}; and Discussion Section~\ref{sec:disc_error}).

    \item \textbf{Develop tools and a framework for neuromechanical control} of the humanoid robot by converting its OpenSim musculoskeletal model to MuJoCo, implementing an SNS-Toolbox two-level spinal CPG with proprioceptive and heel/toe contact feedback, and establishing a staged tuning and verification workflow (Methods Section~\ref{sec:preliminary_sim_methods}, Results Section~\ref{sec:preliminary_sim_results}, and Discussion Section~\ref{sec:synthesis}).

\end{enumerate}

\section{Dissertation Organization}\label{sec:organization}

\begin{reviewchange}
Chapter~\ref{ch:background} establishes the technical and scientific context for the dissertation. Sections~\ref{sec:bg_robots}--\ref{sec:bg_bpamodels} review biomimetic and musculoskeletal robots, artificial-muscle technologies, BPAs, and models of BPA force--length--pressure behavior. Section~\ref{sec:bg_biomechanics} introduces the human lower-limb biomechanics and OpenSim models used as design benchmarks. Sections~\ref{sec:bg_neural}--\ref{sec:bg_tools} then describe spinal central-pattern-generator architectures, sensory-feedback pathways, the Sensory Afferent Database, synthetic nervous systems, and the OpenSim-to-MuJoCo simulation tools used for neuromechanical control. Section~\ref{sec:bg_ethics} closes the chapter by explaining how robotic systems support repeatable neuromechanical experiments that would be impractical or unethical with living subjects.

Chapter~\ref{ch:methods} describes the experimental, analytical, and computational methods that connect actuator behavior to joint performance and neuromechanical control. Sections~\ref{sec:force_model} and~\ref{sec:force_experiment} present the BPA force model and the isometric force-characterization experiment. Sections~\ref{sec:robot_model}--\ref{sec:optimization} describe the pinned-joint and biomimetic-knee test stands, the isometric torque measurements, the hybrid calculation method, the correction terms for constant length offsets, series compliance, and actuator wrapping, and the algorithms used to identify the corrections and redesign actuator routes. Sections~\ref{sec:prior_walker} and~\ref{sec:preliminary_sim_methods} document the earlier AnimatLab walker and the OpenSim-to-MuJoCo and SNS-Toolbox workflows used to construct the neuromechanical models.

Chapter~\ref{ch:results} reports the experimental and computational results. Section~\ref{sec:results_force} presents the BPA force-characterization data and evaluates the resulting maximum-force and nondimensional force models against manufacturer and published predictions. Section~\ref{sec:results_torque} reports the isometric knee-torque measurements and evaluates the correction terms on pinned and biomimetic knees. Section~\ref{sec:preliminary_sim_results} records the instrumented neuromechanical simulations, and Section~\ref{sec:ongoing} reports the full-archive identification and actuator-route redesign. Together, these sections connect the measured actuator behavior, the corrected joint model, the human torque comparison, and the simulation plant used by the controller.

Chapter~\ref{ch:discussion} interprets the results across actuator, transmission, joint, and controller scales. Section~\ref{sec:disc_effects} examines the effect of resting length; Section~\ref{sec:disc_error} evaluates the physical sources of torque-prediction error; and Section~\ref{sec:synthesis} integrates the force-characterization, knee-torque, sensory-database, and neuromechanical-modeling contributions into an actuator-to-controller design methodology. The chapter distinguishes validated conclusions from limitations and identifies which aspects of the framework are transferable to other BPA-driven robots.

Chapter~\ref{ch:futurework} presents concrete research opportunities enabled by the models, experiments, and software tools developed here. Sections~\ref{sec:fw_pipeline} and~\ref{sec:fw_cpg} describe how researchers can reuse the conversion, coupling, and spinal-network tools to compare musculoskeletal plants, sensory pathways, neural lesions, stimulation protocols, and control architectures. Sections~\ref{sec:fw_senses} and~\ref{sec:fw_hardware} describe extensions involving vestibular, ocular, and cerebellar contributions and deployment on embedded hardware. Section~\ref{sec:fw_kickflip} extends the same tools to gait transitions, reactive balance, and dynamic whole-body maneuvers. These are modular applications of completed dissertation contributions, not unfinished dissertation requirements.

Chapter~\ref{ch:conclusion} consolidates the principal contributions: the improved BPA force models, experimentally corrected prediction of isometric knee torque, evaluation of the biomimetic knee and optimized actuator routing, and the tools developed for neuromechanical simulation and control. The chapter emphasizes how the completed experimental and computational foundation supports both improved biomimetic robots and controlled investigations of neuromechanical hypotheses.

Appendix~\ref{app:muscle_calcs} documents BPA force implementation, muscle-path and moment-arm calculations, compliance, and wrapping loss. Appendix~\ref{app:neuron_equations} gives the neuron and synapse equations used by the synthetic-nervous-system tools. Appendix~\ref{app:model_geometry} records the static BPA pathways, joint transformation matrices, bracket frames, stiffness and compliance parameters, and adopted correction-term values.
\end{reviewchange}
